\question[4]
Gegeben ist der folgende Java-Code:

\begin{lstlisting}[language=Java]
int punkte;                  // Zeile 1
punkte = 10;                 // Zeile 2
double schnitt = 2.5;        // Zeile 3
punkte = punkte * 2;         // Zeile 4
punkte = "zwanzig";          // Zeile 5
\end{lstlisting}

\begin{parts}
\part Gib für die Zeilen 1 bis 4 an, ob eine Deklaration, eine Initialisierung oder eine Wertzuweisung vorliegt
(auch mehreres ist möglich).
\part Zeile 5 führt zu einem Fehler. Erläutere, warum.
\part In Python wäre die Wertzuweisung \texttt{punkte = 'zwanzig'} erlaubt. Erläutere den Unterschied zwischen Java und Python.
\end{parts}

\begin{solutionbox}{6cm}
\begin{parts}
\part Zeile 1: Deklaration (Bezeichner \texttt{punkte}, Datentyp \texttt{int}). Zeile 2: Initialisierung.
Zeile 3: Deklaration und Initialisierung in einem Schritt. Zeile 4: Wertzuweisung.
\part \texttt{punkte} ist als \texttt{int} deklariert. In Java ist der Datentyp einer Variable fest, ein String passt nicht in eine
\texttt{int}-Variable. Der Fehler wird schon beim Übersetzen gemeldet.
\part In Java gehört der Datentyp zur Variable und wird bei der Deklaration festgelegt (statische Typisierung).
In Python gibt es keine eigene Deklaration, die Variable entsteht bei der ersten Wertzuweisung. Der Datentyp gehört
zum Wert, deshalb kann eine Variable später einen Wert eines anderen Typs bekommen (dynamische Typisierung).
\end{parts}
\end{solutionbox}
